\section{Profiler：构建可查询的性能与功耗模型}
\label{sec:profiler}
Profiler 的任务是回答：给定模型、TP 布局、频率和请求形状，一个实例能以多大代价承担 P、D 或 M 负载。不同 TP 的通信量、权重占用和 KV 容量不同，因此每个拓扑使用独立 profile；频率也只在相应测量域内比较。模型版本、引擎、硬件和测量文件共同确定 profile 身份，避免规划器混用来源不同的性能数据。

\para{采集维度与执行阶段。}
P 阶段以输入长度 $L$ 和 GPU 频率 $f$ 为主要维度；D 阶段进一步考虑 batch 大小 $B$ 与上下文长度 $C$，逐步解码并记录实际执行形状。M 阶段让 prefill 与 decode 共存，以观测长 prompt 对已有生成请求造成的停顿。采集先完成加载、预热和时钟设置，再分别记录时间与功率窗口，并保留原始记录和独立验证样本。除计算阶段外，还采集空闲、停车、关闭状态的功率，以及唤醒、调频、KV 容量与 P/D 交接信息。

\para{从测量到代价函数。}
记 $T_p(L,f)$ 为 prefill 时间，$T_d(B,C,f)$ 为 decode 单步时间。当前原生 timing 模型按频率进行非负拟合，可简写为：
\begin{align}
 T_p(L,f)&=a_f+b_fL+c_fL^2,\\
 T_d(B,C,f)&=u_f+v_fB+w_fBC.
\end{align}
其中线性项刻画主要工作量，$BC$ 反映上下文相关访存；代码内部的长度归一化已吸收入系数。历史兼容模型还支持 $B^2$ 项和分段曲线，实际查询由绑定的 profile 版本决定。功率同样按频率及 batch/context 建模，M 的功率在具备证据时直接采用混合服务窗口模型。

时间与功率必须对应同一观测区间。对长度为 $H$ 的服务窗口，其平均功率 $\bar P$ 与能量满足
\begin{equation}
 E_{\mathrm{window}}=\int_0^H P(t)\,\mathrm{d}t=H\bar P.
\end{equation}
原生 request-cycle 平均功率覆盖请求调度、计算及间隙，不能直接乘仅包含 CUDA kernel 的 prefill 时间。相应地，整池能耗由匹配的周期或布局窗口估计；停车实例仍计静态功率。这样可以避免把不同口径的测量拼成虚假的节能收益。

\para{传输、验证与查询。}
基础 KV 模型采用固定开销加数据量除以带宽：$T_x(L)=x_0+Lk/v$，其中 $k$ 是每 token 的 KV 字节数，$v$ 是有效带宽。较新的原生组件也可按已测长度插值。然而，M/PD 路径的延迟差包含连接器与调度影响，并非纯物理拷贝时间；第\ref{sec:online}节进一步使用 P 首 token 到 D 首 token 的端点间隔描述真实交接。

拟合与参数选择使用 calibration/tuning 数据，独立 holdout 用于检查预测误差与已选配置。发布时同时绑定允许的频率、长度、batch/context 和拓扑范围；有界版本的越界查询返回缺少覆盖，不能凭插值或理论带宽扩大适用域。规划器还检查静态状态、容量、传输和频率切换组件是否齐备。Profiler 因而提供的不只是一个预测值，而是带有身份、测量口径和有效范围的决策依据；决策记录绑定实际使用的 profile，便于回查。
